\begin{table}[H]
\centering
\caption{\textbf{Propensity-model candidates: common support, effective sample size, balance, and weights.}
Candidates were compared without using outcomes; the model with the smaller maximum weighted absolute
standardized mean difference was selected, with the 63-leaf model preferred when the difference was no
more than 0.005. Support is the share of eligible papers with an estimated probability of narrower-scope
publication within 0.05--0.95. Weights are inverse-probability weights within support. The
reference-adjusted model adds the classified final reference list to the adjustment set and defines its
own support; the regenerated fit repeats the primary specification a second time to store
paper-level scores.}
\label{tab:propensity}
\footnotesize
\setlength{\tabcolsep}{4pt}
\begin{tabular}{@{}lrrrrrrrr@{}}
\toprule
 & & & \multicolumn{2}{c}{Effective sample size} & Max. weighted & \multicolumn{3}{c}{Weight percentile} \\
\cmidrule(lr){4-5}\cmidrule(l){7-9}
Propensity model & Support, \% & Papers & Broader & Narrower & abs. SMD & 50th & 95th & 99th \\
\midrule
63 leaves (selected; primary) & 50.1 & 3,818,173 & 671,512 & 1,040,789 & 0.180 & 1.33 & 7.58 & 15.29 \\
255 leaves & 37.2 & 2,832,802 & 500,631 & 694,407 & 0.268 & 1.39 & 10.29 & 17.17 \\
63 leaves, reference-adjusted & 49.7 & 3,782,942 & 668,873 & 1,006,812 & 0.176 & 1.31 & 7.63 & 15.34 \\
63 leaves, regenerated fit & 50.2 & 3,827,491 & 670,756 & 1,021,386 & 0.194 & 1.32 & 7.69 & 15.43 \\
\bottomrule
\end{tabular}
\end{table}
